\documentclass[../../main2.tex]{subfiles}

\begin{document}

\chapter{无你检验：如果没有你，还会发生吗}
\label{ch:butfor}

\begin{motivation}
	Part II 遗留的问题：因果的手艺拿到了，但「各占多少」还是一句嘴上话。科学界可以含糊，有一个地方躲不掉——法庭。判决书不许写「大概、也许、都有点责任」，必须写「七成是他的责任」。\\
	\textbf{核心动机}：看人类怎么被「必须给数字」逼了几千年，逼出一套\textbf{无你检验}（but-for test）——它就是「各占多少」这门手艺最早成熟的形态。本书的公式不是发明，是抄法庭的作业。
\end{motivation}

\begin{logicmap}
\begin{tikzpicture}[node distance=0.85cm, every node/.style={draw, rectangle, rounded corners, align=center, font=\small}]
\node (q) {「差不多」不行：必须给数字};
\node (c) [below=of q] {无你检验：拿掉你，还发生吗};
\node (d) [below=of c] {同一动作的公式：条件概率差};
\node (p) [below=of d] {边界：无你检验不够用时};
\node (n) [below=of p] {一堆原因怎么分账（抛给第 9 章）};
\draw[-{Stealth}] (q) -- (c);
\draw[-{Stealth}] (c) -- (d);
\draw[-{Stealth}] (d) -- (p);
\draw[-{Stealth}] (p) -- (n);
\end{tikzpicture}
\end{logicmap}

\section{当「差不多」不行的时候}

\begin{readingnote}
	你有没有想过，为什么「不得杀人」四个字人人会背，真要定责时却要几百页卷宗？\\
	因为背规则只花一秒，\textbf{分责任}要花掉整个审判。
\end{readingnote}

先看一个真实压力下的场景。两车相撞，一辆逆行，一辆超速，骑电动车的第三方受伤。三方都有过错。法官必须写下的不是「都有责任」，而是：逆行方承担六成，超速方三成，电动车一成。

为什么必须给数字？因为\textbf{纠纷必须了结}。赔偿款是一个确定的数，谁掏多少是一个确定的比例。「差不多」在法庭上不是答案，是渎职。这个压力扛了上千年——科学界直到近现代才有的「必须可检验」，法律界从汉谟拉比法典起就在扛。

\begin{tcolorbox}[colback=green!5, title=\textit{生活类比}]
	科学论文可以说「显著相关」，像体检报告写「建议复查」；判决书必须写「七成责任、赔偿三十八万」，像医生上手术台——刀口开多深，不能「建议一下」。压力不同，手艺精度不同。
\end{tcolorbox}

三个不同人群的「必须给数字」现场，今天还在发生：

\begin{itemize}
	\item \textbf{上班族}：项目黄了，复盘会上老板问「谁的责任、几成」——你躲不过去。
	\item \textbf{学生}：小组作业划水，分数怎么分——「都干了」不是答案。
	\item \textbf{结构内行}：事故调查报告的最后一栏永远是「主要原因 / 次要原因 / 各占比例」。
\end{itemize}

\paragraph*{逻辑衔接}
	压力有了：必须给数字。那人类被逼出了什么动作？下一节看那把千年镊子。

\section{无你检验：拿掉你，还发生吗}

\begin{readingnote}
	你有没有想过，为什么法庭问的是「如果他没有开枪，死者还会死吗」，而不是「他是不是坏人」？\\
	因为责任不问善恶，问\textbf{反事实}。
\end{readingnote}

法律给的动作，简单到不像千年智慧：

\begin{quote}
	\textbf{无你检验}（but-for test）：如果没有被告的行为，损害还会发生吗？\\
	不会——行为是原因；还会——行为不是（这个损害的）原因。
\end{quote}

这就是第 6 章的「平行宇宙对比」，只不过法庭每周开庭都在做。几个案例走一遍：

\begin{itemize}
	\item \textbf{工伤}：脚手架没绑安全绳，工人坠落。「没这根绳，他会摔吗？」——会摔（绳不承重），那绳不是原因？不对：本案里他是\textbf{从没绑绳的缺口}掉下去的——没有这次违规，他不会在这个位置摔。\textbf{过}。
	\item \textbf{医疗}：医生误诊拖延三周，患者病情恶化。「不拖延，病能治好吗？」——早期治愈率高，拖延改变了概率。\textbf{有份}。
	\item \textbf{投毒}：甲投了致死量毒药，乙也投了致死量毒药，人死了。「拿掉甲，人还会死吗？」——会，因为乙的量也够。无你检验在这里会说「甲无责」——荒唐！这个洞，是下一章「多因分账」的入口。
	\item \textbf{古}：唐代的「保辜」制度——打人后伤者若干日内死，才算殴杀。等一等再判，就是在等「没有这一拳，他会死吗」的证据。
\end{itemize}

\begin{questionbox}
	\textit{现在你可能想反驳：「法官定责是拍脑袋，哪真算概率？」}\\
	\textbf{对，判决书里没有公式。}但看推理的\textbf{结构}：「若无被告行为，损害是否仍会发生」——问的正是两个世界的对比。法官不写概率，不代表这个动作不是概率动作；代表它的精度被压缩成了「可能 / 不可能」两档。本书要做的是把这个压缩档位还原。
\end{questionbox}

\paragraph*{逻辑衔接}
	动作有了：对比有你没你两个世界。怎么把它变成能算的量？下一节——全书第一个公式，配一个能算出数的案例。

\section{同一动作，写成公式}
\label{sec:butfor-formula}

先看案例，再看公式（本书铁律：公式前面必须站着真实案例）。

\begin{examplebox}
	\textbf{案例：工厂排污与鱼塘}。鉴定结论：工厂排污期间，鱼苗死亡概率约三成；对照未排污期，死亡概率约百分之五。问：排污「占多少」？
\end{examplebox}

现在把「拿掉排污」这个动作写成一个数：

\begin{mathbox}{贡献 = 条件概率差（无你检验的公式形态）}
	\[
	C_A \;=\; P(D \mid A) \;-\; P(D \mid \neg A)
	\]
	读法：$A$ = 行为（排污），$D$ = 损害（鱼苗死亡）。
	$P(D \mid A)$ = 有行为时损害多大概率；$P(D \mid \neg A)$ = 平行宇宙里没行为时多大概率。\\[0.5em]
	\textbf{代入案例}：$C_A = 0.30 - 0.05 = 0.25$。\\
	解释：排污把死亡概率\textbf{抬高了 25 个百分点}——这 0.25 就是排污的「份额」。没排污的基准线 0.05 不算它的账。\\[0.5em]
	\textbf{为什么这样定义好？}它把「我觉得他有责任」钉成了可对比的量：两个工厂、两起事故，谁的 $C$ 大谁的责任重，可以排队了。「责任」从感受变成了数。
\end{mathbox}

\mathlink{app:contrib}{严格定义、假设与算例的完整版见数学附录 A.4。}

再走一个更日常的：

\begin{examplebox}
	\textbf{案例：迟到与奖金}。你迟到两小时，错过关键汇报，丢掉当月奖金。不迟到的平行宇宙里，你大概率能拿奖。$P(\text{丢单} \mid \text{迟到}) \approx 0.8$，$P(\text{丢单} \mid \text{不迟到}) \approx 0.1$，则 $C_{\text{迟到}} = 0.8 - 0.1 = 0.7$。七成锅，你自己背。
\end{examplebox}

注意这个定义的两个脾气：

\begin{itemize}
	\item \textbf{它只认「可拿掉」}：拿掉 A 世界就不同，A 才有账。天生的东西（遗传病史）没法拿掉，没法进公式。
	\item \textbf{它是概率差，不是道德分}：$C$ 大不等于「坏」，只等于「起的作用大」。台风天的雨也有 $C$，雨不缺德。
\end{itemize}

\begin{questionbox}
	\textit{现在你可能想反驳：「两个世界的概率，现实里只有一个世界啊，怎么观测？」}\\
	\textbf{问到点子上了。}平行宇宙观测不到——所以 $C$ 永远是\textbf{估}出来的。第 7 章三把镊子（自然实验、剂量反应、机制）就是估 $C$ 的工具箱；第 9 章会告诉你，多个原因同时在场时，估计还要更小心。
\end{questionbox}

\paragraph*{逻辑衔接}
	「各占多少」第一次变成了能算的量。但刚才投毒案里无你检验出了洋相：两个凶手各自「拿掉都还死」，难道都无责？下一章正面拆这个洞。

\section{边界：无你检验什么时候失灵}

\begin{readingnote}
	你有没有想过，一个判了几千年的检验，会有什么\textbf{系统性}的漏洞？\\
	漏洞恰恰出在它最得意的地方——「拿掉你，还发生吗」。
\end{readingnote}

三类失灵，心里先记下（下一章逐个拆）：

\begin{itemize}
	\item \textbf{双充分原因}（overdetermination）：两把独立的火同时点燃同一栋楼——拿掉哪把，楼都烧。两把火都无责？显然不对。
	\item \textbf{合力原因}：单独不下毒、混合才致死的两种药——拿掉谁都无害，合起来要命。谁的锅？
	\item \textbf{背景原因}：地心引力「导致」了摔伤。拿掉引力人不摔——难道判引力有责？\textbf{因果需要背景，责任不需要}——法律靠「可责性」「近因」把背景挡在门外。
\end{itemize}

\lowfreq{诚实标注：法律的「近因」「相当因果关系」等概念，就是给无你检验打的补丁。补丁列表历来有争议——本书只取其逻辑内核，不装作它完美。}

\paragraph*{逻辑衔接}
	无你检验管用，但只有单人剧本。现实的事故、革命、爆款，从来是多因同台。下一章：多因分账——交互项、分摊方案，全书最硬的一块。

\begin{thinkbox}
\begin{enumerate}
	\item 用无你检验拆一次你自己的「事故」：上个月某件不顺的事（迟到、丢单、吵架），列出你怀疑的「行为」，逐个问「拿掉它，损失还会发生吗」。哪个是真原因，哪个只是路过？
	\item \textbf{反事实}：如果唐代没有「保辜」制度——打人当场定罪、不等伤势发展——法医学会提前出现吗？「无你检验」这门手艺会更糙还是更精细？
	\item \textbf{反事实}：如果工厂排污案里，鉴定给出的不是「三成对百分之五」，而是「排污期和对照期死亡率一样」，$C_A$ 是多少？还判得下去吗？
\end{enumerate}
\end{thinkbox}

\end{document}
